\documentclass[addpoints]{exam}
\usepackage[spanish]{babel}
\begin{document}
	\begin{questions}
		\question Given the equation \(x^n + y^n = z^n\) for \(x,y,z\) and \(n\) positive
		integers. 
		\begin{parts}
			\part[5] For what values of \(n\) is the statement in the previous question true?
			\vspace{\stretch{1}}
			
			\part[2 \half] For \(n=2\) there's a theorem with a special name. What's that name?
			\vspace{\stretch{1}}
			
			\part[2 \half] What famous mathematician had an elegant proof for this theorem but there was
			not enough space in the margin to write it down?
			\vspace{\stretch{1}}
			
		\end{parts}
		
		\droptotalpoints
		
		\question[20] Compute \[\int_{0}^{\infty} \frac{\sin(x)}{x}\]
	\end{questions}
	\begin{questions}
		
		\question Given the equation \(x^n + y^n = z^n\) for \(x,y,z\) and \(n\) positive
		integers. 
		\begin{parts}
			\part[5] For what values of \(n\) is the statement in the previous question true?
			\vspace{\stretch{1}}
			
			\part[2 \half] For \(n=2\) there's a theorem with a special name. What's that name?
			\vspace{\stretch{1}}
			
			\bonuspart[2 \half] What famous mathematician had an elegant proof for this theorem but there was
			not enough space in the margin to write it down?
			\vspace{\stretch{1}}
			
		\end{parts}
		
		\droptotalpoints
		
		\question[20] Compute \[\int_{0}^{\infty} \frac{\sin(x)}{x}\]
		
		\vspace{\stretch{1}}
		
		\bonusquestion[30] Prove that the real part of all non-trivial zeros of the function 
		\(\zeta(z)\) is \(\frac{1}{2}\)
		\vspace{\stretch{1}}
		
	\end{questions}
	\begin{center}
		\combinedgradetable[h][questions]
	\end{center}
	
	
	
Here we'll use \verb|\pointsinmargin|
\begin{questions}
	\pointsinmargin
	\question Given the equation \(x^n + y^n = z^n\) for \(x,y,z\) and \(n\) positive
	integers. 
	\begin{parts}
		\part[5] For what values of \(n\) is the statement in the previous question true?
		\vspace{40pt}
		
		\part[2 \half] For \(n=2\) there's a theorem with a special name. What's that name?
		\vspace{40pt}
		
		\part[2 \half] What famous mathematician had an elegant proof for this theorem but there was
		not enough space in the margin to write it down?
		
		\vspace{40pt}
	\end{parts}
	
	Now we'll use \verb|\pointsinrightmargin|
	
	\pointsinrightmargin
	\question Given the equation \(x^n + y^n = z^n\) for \(x,y,z\) and \(n\) positive
	integers. 
	\begin{parts}
		\part[5] For what values of \(n\) is the statement in the previous question true?
		\vspace{40pt}
		
		\part[2 \half] For \(n=2\) there's a theorem with a special name. What's that name?
		\vspace{40pt}
		
		\part[2 \half] What famous mathematician had an elegant proof for this theorem but there was
		not enough space in the margin to write it down?
		integers. 
	\end{parts}
\end{questions}
\end{document}